\documentclass[12pt,a4paper]{article}
\usepackage{bendingforks}

% =====================================================================
%  STUDIO DI FATTIBILITÀ - documentazione interna
%  I testi in rosso (\ph{...}) sono segnaposto: vanno sostituiti tutti
%  prima della consegna. Le date vanno sempre nel formato AAAA-MM-GG.
% =====================================================================
\bfsetup{Studio di Fattibilità}{Analisi dei capitolati e motivazione della scelta}{0.1.0}{Documento}{Interno}
\bfredattori{\ph{Nome Cognome} \\ & \ph{Nome Cognome}}
\bfverificatori{\ph{Nome Cognome}}
\bfdestinatari{\emph{Gruppo Bending Forks} \\ & \emph{Prof. Tullio Vardanega} \\ & \emph{Prof. Riccardo Cardin}}
\bfabstract{Analisi dei capitolati proposti per il progetto didattico di Ingegneria del Software A.A.~2026/27: per ciascuno si valutano obiettivi, tecnologie, vincoli, rischi e fattibilità, e si motivano la scelta del capitolato C3 \emph{Fossdot} e lo scarto degli altri.}

\begin{document}
\bfcopertina
\bfindice

% ---------- Registro delle versioni ----------
\begin{bfregistro}
0.1.0 & \ph{2026-10-27} & \ph{Nome Cognome} & Creazione della struttura del documento \bfsep
\end{bfregistro}

% =====================================================================
\section{Introduzione}

\subsection{Scopo del documento}
Il presente documento raccoglie l'analisi dei capitolati d'appalto proposti per il progetto didattico del corso di Ingegneria del Software (A.A.~2026/27). Per ciascun capitolato si individuano obiettivi, tecnologie richieste, vincoli, aspetti positivi, criticità e rischi, valutandone la fattibilità tecnica ed economica rispetto alle competenze e alle risorse del gruppo \emph{Bending Forks}. Il documento motiva infine la scelta del capitolato per cui il gruppo presenta la propria candidatura e le ragioni dello scarto degli altri.

\subsection{Glossario}
Per evitare ambiguità, i termini tecnici usati nel documento sono definiti nel \emph{Glossario}. \ph{Indicare come vengono marcati i termini presenti nel Glossario, ad es. con una G a pedice, oppure togliere questa sezione se il Glossario non è ancora previsto.}

\subsection{Elenco dei capitolati analizzati}

\begin{center}
\begin{tabularx}{\linewidth}{>{\hsize=0.35\hsize\centering\arraybackslash}X>{\hsize=1.25\hsize\raggedright\arraybackslash}X>{\hsize=1.0\hsize\raggedright\arraybackslash}X>{\hsize=0.4\hsize\centering\arraybackslash}X}
\bfheadrow\bfh{Codice} & \bfh{Capitolato} & \bfh{Proponente} & \bfh{Colloquio} \bfsep
C0 & AIRA: AI Requests Automation & Sanmarco Informatica & -- \bfsep
C1 & AURA & Vimar & -- \bfsep
C2 & Concierge (NEXUM) & Eggon & -- \bfsep
C3 & \textbf{Fossdot} & \textbf{Uranion} & \textbf{2026-10-21} \bfsep
C4 & Gestione Veicoli e Servizi Post-Vendita & M31 & -- \bfsep
C5 & RareCare & Miriade & -- \bfsep
C6 & Caching semantico e dashboard FinOps & Var Group & -- \bfsep
C7 & Smart Dashboard & Ergon Informatica & 2026-10-22 \bfsep
C8 & Sonar & Halue & -- \bfsep
C9 & SpecCheck & Zucchetti & -- \bfsep
\end{tabularx}
\end{center}

\subsection{Riferimenti}
\begin{itemize}
\tightlist
\item Regolamento del progetto didattico: \ph{URL}
\item Capitolati d'appalto A.A.~2026/27: \ph{URL della pagina dei capitolati sul sito del corso}
\item Verbali degli incontri con le aziende: Appendice~\ref{app:verbali}
\end{itemize}

% =====================================================================
\section{Analisi del capitolato C3 -- \emph{Fossdot}}
\label{sec:c3}

\subsection{Descrizione generale}
\textbf{Proponente:} \emph{Uranion} \\
\ph{Descrizione in 5-10 righe del prodotto richiesto e del contesto in cui si inserisce.}

\subsection{Obiettivi e finalità}
\ph{Cosa deve ottenere il prodotto (MVP e versione completa) e per chi.}

\subsection{Tecnologie richieste}
\begin{itemize}
\tightlist
\item \ph{Tecnologia} -- \ph{a cosa serve, se è imposta o solo suggerita}
\item \ph{Tecnologia} -- \ph{...}
\end{itemize}

\subsection{Vincoli e requisiti}
\begin{itemize}
\tightlist
\item \textbf{Requisiti obbligatori:} \ph{elenco sintetico}
\item \textbf{Requisiti opzionali:} \ph{elenco sintetico}
\item \textbf{Vincoli tecnologici e di processo:} \ph{...}
\end{itemize}

\subsection{Aspetti positivi}
\begin{itemize}
\tightlist
\item \ph{Aspetto positivo}
\end{itemize}

\subsection{Criticità e rischi}

\begin{center}
\begin{tabularx}{\linewidth}{>{\hsize=0.7\hsize\raggedright\arraybackslash}X>{\hsize=1.5\hsize\raggedright\arraybackslash}X>{\hsize=0.45\hsize\centering\arraybackslash}X>{\hsize=0.45\hsize\centering\arraybackslash}X>{\hsize=0.9\hsize\raggedright\arraybackslash}X}
\bfheadrow\bfh{Tipo} & \bfh{Rischio} & \bfh{Prob.} & \bfh{Impatto} & \bfh{Mitigazione} \bfsep
Tecnologico & \ph{Descrizione} & \ph{A/M/B} & \ph{A/M/B} & \ph{Strategia} \bfsep
Personale & \ph{Descrizione} & \ph{A/M/B} & \ph{A/M/B} & \ph{Strategia} \bfsep
Requisiti & \ph{Descrizione} & \ph{A/M/B} & \ph{A/M/B} & \ph{Strategia} \bfsep
Tempistiche & \ph{Descrizione} & \ph{A/M/B} & \ph{A/M/B} & \ph{Strategia} \bfsep
\end{tabularx}
\end{center}

\subsection{Valutazione della fattibilità}
\subsubsection{Competenze disponibili}
\ph{Cosa il gruppo sa già fare (linguaggi, strumenti, esperienze) e chi ha quelle competenze.}
\subsubsection{Competenze da acquisire}
\ph{Gap analysis: cosa serve imparare e come (documentazione, formazione dell'azienda, prove).}
\subsubsection{Stima dei tempi di apprendimento}
\ph{Ore stimate per argomento e per persona.}
\subsubsection{Piattaforme e strumenti}
\ph{Idoneità di hardware, servizi e strumenti all'esecuzione del progetto; cosa fornisce l'azienda.}

\subsection{Valutazione economica preliminare}
\ph{Rapporto tra ore necessarie e complessità; stima di massima del costo e confronto con il vincolo di 12.000\,€.}

\subsection{Conclusioni}
\ph{Giudizio sintetico sul capitolato.}

% =====================================================================
\section{Analisi del capitolato C7 -- \emph{Smart Dashboard}}
\label{sec:c7}

\subsection{Descrizione generale}
\textbf{Proponente:} \emph{Ergon Informatica} \\
\ph{Descrizione in 5-10 righe del prodotto richiesto e del contesto in cui si inserisce.}

\subsection{Obiettivi e finalità}
\ph{Cosa deve ottenere il prodotto e per chi.}

\subsection{Tecnologie richieste}
\begin{itemize}
\tightlist
\item \ph{Tecnologia} -- \ph{a cosa serve, se è imposta o solo suggerita}
\end{itemize}

\subsection{Vincoli e requisiti}
\ph{Requisiti obbligatori, opzionali e vincoli.}

\subsection{Aspetti positivi}
\ph{...}

\subsection{Criticità e rischi}
\ph{Stessa tabella della sezione~\ref{sec:c3}, se utile.}

\subsection{Valutazione della fattibilità}
\ph{Competenze disponibili, competenze da acquisire, stima dei tempi di apprendimento.}

\subsection{Valutazione economica preliminare}
\ph{...}

\subsection{Conclusioni}
\ph{...}

% ---------------------------------------------------------------------
%  Per ogni altro capitolato analizzato copiare la sezione C7 qui sopra
%  e cambiare codice, nome e proponente.
% ---------------------------------------------------------------------

% =====================================================================
\section{Capitolato scelto: C3 -- \emph{Fossdot}}

\subsection{Motivazioni della scelta}
\begin{itemize}
\item \textbf{Allineamento con le competenze del gruppo:} \ph{...}
\item \textbf{Interesse tecnico e professionale:} \ph{...}
\item \textbf{Rapporto tra complessità e fattibilità:} \ph{...}
\item \textbf{Qualità del rapporto con il proponente:} \ph{...}
\end{itemize}

\subsection{Strategie di mitigazione dei rischi}
\ph{Per i rischi principali del capitolato: come vengono prevenuti e qual è il piano B se l'approccio principale fallisce.}

\subsection{Piano di formazione preliminare}

\begin{center}
\begin{tabularx}{\linewidth}{>{\hsize=1.3\hsize\raggedright\arraybackslash}X>{\hsize=1.0\hsize\raggedright\arraybackslash}X>{\hsize=0.9\hsize\raggedright\arraybackslash}X>{\hsize=0.8\hsize\centering\arraybackslash}X}
\bfheadrow\bfh{Argomento} & \bfh{Modalità} & \bfh{Chi} & \bfh{Entro} \bfsep
\ph{Argomento} & \ph{Studio individuale / workshop azienda} & \ph{Nomi} & \ph{AAAA-MM-GG} \bfsep
\ph{Argomento} & \ph{...} & \ph{Nomi} & \ph{AAAA-MM-GG} \bfsep
\end{tabularx}
\end{center}

\subsection{Impegno orario stimato}
\ph{Ore produttive previste per persona (tra 80 e 95) e totali; rimando al documento \emph{Preventivo dei costi e assunzione degli impegni}.}

% =====================================================================
\section{Capitolati scartati}

\subsection{C7 -- \emph{Smart Dashboard}}
\ph{Motivi dello scarto.}
\subsection{C0 -- \emph{AIRA}}
\ph{Motivi dello scarto.}
\subsection{C1 -- \emph{AURA}}
\ph{Motivi dello scarto.}
\subsection{C2 -- \emph{Concierge}}
\ph{Motivi dello scarto.}
\subsection{C4 -- \emph{Gestione Veicoli e Servizi Post-Vendita}}
\ph{Motivi dello scarto.}
\subsection{C5 -- \emph{RareCare}}
\ph{Motivi dello scarto.}
\subsection{C6 -- \emph{Caching semantico e dashboard FinOps}}
\ph{Motivi dello scarto.}
\subsection{C8 -- \emph{Sonar}}
\ph{Motivi dello scarto.}
\subsection{C9 -- \emph{SpecCheck}}
\ph{Motivi dello scarto.}

% =====================================================================
\appendix
\section{Contatti con le aziende}

\begin{center}
\begin{tabularx}{\linewidth}{>{\hsize=0.8\hsize\raggedright\arraybackslash}X>{\hsize=1.0\hsize\raggedright\arraybackslash}X>{\hsize=0.6\hsize\centering\arraybackslash}X>{\hsize=0.9\hsize\raggedright\arraybackslash}X>{\hsize=0.7\hsize\raggedright\arraybackslash}X}
\bfheadrow\bfh{Azienda} & \bfh{Referente} & \bfh{Data} & \bfh{Modalità} & \bfh{Capitolato} \bfsep
Uranion & Giacomo Cariello & 2026-10-21 & Google Meet & C3 \bfsep
Ergon Informatica & Gianluca Carlesso & 2026-10-22 & Microsoft Teams & C7 \bfsep
\end{tabularx}
\end{center}

\section{Verbali dei colloqui esplorativi}
\label{app:verbali}
\begin{itemize}
\tightlist
\item Verbale Esterno N.~1 -- Uranion, 2026-10-21: \ph{link al PDF nella repository}
\item Verbale Esterno N.~2 -- Ergon Informatica, 2026-10-22: \ph{link al PDF nella repository}
\end{itemize}

\end{document}
